\documentclass[12pt]{article}
\usepackage{amsmath, amssymb}
\usepackage[margin=1in]{geometry}
\setlength{\parskip}{6pt}
\title{QUBO Formulation: Partition Problem}
\date{}
\begin{document}
\maketitle

\subsection*{Partition Problem}

\paragraph{Problem Statement.}
Given a collection of $N$ positive integers $\{a_0, a_1, \dots, a_{N-1}\}$, the goal is to partition the index set $\{0, \dots, N-1\}$ into two disjoint subsets $\mathcal{S}_0$ and $\mathcal{S}_1$ such that the absolute difference between the sum of the elements in $\mathcal{S}_1$ and the sum of the elements in $\mathcal{S}_0$ is minimized.

\paragraph{Decision Variables.}
For each index $i \in \{0, \dots, N-1\}$, we introduce a binary decision variable $x_i \in \{0, 1\}$. The variable $x_i$ encodes the assignment of the $i$-th integer $a_i$: $x_i = 1$ indicates that $a_i$ is assigned to subset $\mathcal{S}_1$, while $x_i = 0$ indicates assignment to subset $\mathcal{S}_0$.

\paragraph{QUBO Derivation.}
Step 1 --- The sum of elements in $\mathcal{S}_1$ is $\sum_{i=0}^{N-1} a_i x_i$ and in $\mathcal{S}_0$ is $\sum_{i=0}^{N-1} a_i (1 - x_i)$. Their difference is $\sum_{i=0}^{N-1} a_i (2x_i - 1)$. Minimizing the absolute difference is equivalent to minimizing its square:
\begin{align}
    \min_{x \in \{0,1\}^N} \left( \sum_{i=0}^{N-1} a_i (2x_i - 1) \right)^2.
\end{align}

Step 2 --- The problem imposes no structural constraints on the assignment of elements to subsets.

Step 3 --- We expand the squared objective directly. Let $C = \sum_{k=0}^{N-1} a_k$ denote the total sum of all integers. Expanding the square yields:
\begin{align}
    \left( \sum_{i=0}^{N-1} a_i (2x_i - 1) \right)^2 &= \sum_{i=0}^{N-1} a_i^2 (2x_i - 1)^2 + \sum_{i \neq j} a_i a_j (2x_i - 1)(2x_j - 1).
\end{align}
Using the idempotent property of binary variables, $x_i^2 = x_i$, we simplify $(2x_i - 1)^2 = 4x_i^2 - 4x_i + 1 = 1$. The cross terms expand as $(2x_i - 1)(2x_j - 1) = 4x_i x_j - 2x_i - 2x_j + 1$. Substituting these back:
\begin{align}
    H(x) &= \sum_{i=0}^{N-1} a_i^2 + \sum_{i \neq j} a_i a_j (4x_i x_j - 2x_i - 2x_j + 1) \\
         &= \sum_{i=0}^{N-1} a_i^2 + 4 \sum_{0 \leq i < j \leq N-1} a_i a_j x_i x_j - 2 \sum_{i \neq j} a_i a_j x_i - 2 \sum_{i \neq j} a_i a_j x_j + \sum_{i \neq j} a_i a_j.
\end{align}
Collecting the linear terms for each $x_i$, we note that $x_i$ appears in the cross sums whenever $j \neq i$. Thus, the coefficient of $x_i$ is:
\begin{align}
    -2 \sum_{j \neq i} a_i a_j - 2 \sum_{j \neq i} a_j a_i = -4 a_i \sum_{j \neq i} a_j = -4 a_i (C - a_i) = 4a_i^2 - 4C a_i.
\end{align}
The constant terms combine as $\sum_{i=0}^{N-1} a_i^2 + \sum_{i \neq j} a_i a_j = \left( \sum_{i=0}^{N-1} a_i \right)^2 = C^2$.

Step 4 --- Combining all contributions, the QUBO Hamiltonian is:
\begin{align}
    H(x) = \sum_{i=0}^{N-1} (4a_i^2 - 4C a_i) x_i + \sum_{0 \leq i < j \leq N-1} 4a_i a_j x_i x_j + C^2.
\end{align}

Step 5 --- Reading off the coefficients from the standard QUBO form $H(x) = \sum_{i} Q_{ii} x_i + \sum_{i < j} Q_{ij} x_i x_j + c$, we obtain the general closed-form expressions:
\begin{align}
    Q_{ii} &= 4a_i^2 - 4C a_i, \quad \forall i \in \{0, \dots, N-1\}, \\
    Q_{ij} &= 4a_i a_j, \quad \forall 0 \leq i < j \leq N-1, \\
    c &= C^2.
\end{align}

\paragraph{Penalty Weight Justification.}
The problem formulation contains no hard constraints on the variable assignments. Every binary vector $x \in \{0,1\}^N$ corresponds to a valid partition of the integers into two subsets. Consequently, no penalty terms are required, and the Hamiltonian consists solely of the unconstrained objective expansion. The number of variables is exactly $N$.

\paragraph{Correctness Argument.}
Minimizing $H(x)$ over the Boolean hypercube $\{0,1\}^N$ recovers the optimal partition because $H(x)$ is a strictly increasing function of the absolute difference $|S_1 - S_0|$ for non-negative differences. Since the square function preserves the ordering of non-negative values, the global minimum of $H(x)$ occurs precisely at the assignment $x^*$ that minimizes $|S_1 - S_0|$. All $2^N$ possible assignments are feasible, and the quadratic structure ensures that the energy landscape directly reflects the objective function's values without introducing artificial infeasibilities.

\paragraph{Illustrative Example.}
Consider a concrete instance with $N=3$ integers $\{a_0, a_1, a_2\} = \{2, 3, 4\}$. The total sum is $C = 2 + 3 + 4 = 9$. Applying the general formulas yields the following QUBO coefficients:
\begin{align}
    Q_{00} &= 4(2)^2 - 4(9)(2) = -56, & Q_{11} &= 4(3)^2 - 4(9)(3) = -72, & Q_{22} &= 4(4)^2 - 4(9)(4) = -80, \\
    Q_{01} &= 4(2)(3) = 24, & Q_{02} &= 4(2)(4) = 32, & Q_{12} &= 4(3)(4) = 48.
\end{align}
The constant offset is $c = 9^2 = 81$. The resulting Hamiltonian is:
\begin{align}
    H(x_0, x_1, x_2) = -56x_0 - 72x_1 - 80x_2 + 24x_0x_1 + 32x_0x_2 + 48x_1x_2 + 81.
\end{align}
Enumerating all $2^3 = 8$ binary assignments, the minimum energy is $H(0, 1, 1) = 1$, achieved at $x^* = (0, 1, 1)$. This corresponds to assigning $a_0=2$ to subset $\mathcal{S}_0$ and $\{a_1=3, a_2=4\}$ to subset $\mathcal{S}_1$, yielding subset sums $2$ and $7$ with an absolute difference of $|7-2|=5$, which is the optimal partition for this instance.

\end{document}